\documentclass[11pt,a4paper]{article}

\usepackage[T1]{fontenc}
\usepackage[utf8]{inputenc}
\usepackage[ngerman]{babel}

\usepackage{amsmath,amssymb,amsthm,amsfonts}
\usepackage{mathtools}
\usepackage{geometry}
\usepackage{physics}
\usepackage{bm}
\usepackage{graphicx}
\usepackage{hyperref}
\usepackage{enumitem}

\geometry{margin=2.5cm}

\title{Der EABC-Tensorraum:\\
	Eine spektrale und geometrische Theorie symmetrischer EABC-Kopplungen}

\author{Thomas Hoffbauer}

\date{\today}

\theoremstyle{definition}
\newtheorem{definition}{Definition}
\newtheorem{satz}{Satz}
\newtheorem{lemma}{Lemma}
\newtheorem{korollar}{Korollar}
\newtheorem{bemerkung}{Bemerkung}
\newtheorem{beispiel}{Beispiel}

\begin{document}
	
	\maketitle
	
	\begin{abstract}
		In diesem Artikel wird eine tensorielle Erweiterung des EABC-Formalismus entwickelt. Ausgangspunkt ist ein symmetrischer, spurloser Tensorraum, dessen fundamentale Operatoren aus der EABC-Struktur hervorgehen. Die entstehende Geometrie besitzt eine bemerkenswert einfache spektrale Struktur: Alle Eigenwerte hängen ausschließlich von zwei fundamentalen Invarianten ab. 
		
		Die Theorie verbindet algebraische Eigenschaften der EABC-Struktur mit Tensoranalysis, Quaternionengeometrie, Spektraltheorie und physikalischen Kopplungsmodellen. Insbesondere zeigt sich, dass die kubische Invariante eine orientierte Triple-Kohärenz beschreibt, während die quadratische Invariante als Gesamtenergie bzw. Gesamtkohärenz interpretiert werden kann.
		
		Darüber hinaus wird gezeigt, dass die Diskriminante des Systems stets nichtnegativ ist, wodurch eine vollständig reelle Spektralstruktur garantiert wird.
	\end{abstract}
	
	\tableofcontents
	
	\newpage
	
	\section{Einleitung}
	
	Die EABC-Struktur beschreibt eine algebraische Zerlegung diskreter Zustandsräume in vier fundamentale Klassen:
	
	\[
	E,\quad A,\quad B,\quad C
	\]
	
	mit den Relationen
	
	\[
	A^2=B^2=C^2=E
	\]
	
	sowie
	
	\[
	AB=C,\qquad AC=B,\qquad BC=A.
	\]
	
	Diese Struktur entspricht algebraisch der Klein'schen Vierergruppe und besitzt zugleich tiefe Verbindungen zu Quaternionen, modularen Restklassen und spektralen Operatorräumen.
	
	Ziel dieses Artikels ist die Entwicklung eines kontinuierlichen Tensorformalismus auf Basis der EABC-Operatoren.
	
	Dabei entsteht ein bemerkenswert einfacher, aber äußerst reichhaltiger Tensorraum mit:
	
	\begin{itemize}
		\item symmetrischer Struktur,
		\item spurlosen Kopplungsoperatoren,
		\item exakt lösbarer Spektralgeometrie,
		\item zwei fundamentalen Invarianten,
		\item natürlicher Quaternionenstruktur,
		\item physikalischer Interpretierbarkeit.
	\end{itemize}
	
	\section{Die fundamentale EABC-Struktur}
	
	Wir betrachten die vier Basisoperatoren:
	
	\[
	E,\quad A,\quad B,\quad C.
	\]
	
	Die algebraischen Relationen lauten:
	
	\[
	A^2=B^2=C^2=E
	\]
	
	und
	
	\[
	AB=C,\qquad BC=A,\qquad AC=B.
	\]
	
	Diese Struktur besitzt:
	
	\begin{itemize}
		\item eine involutive Symmetrie,
		\item eine modulare Restklassenstruktur,
		\item eine natürliche Interpretation als diskreter Orientierungsraum.
	\end{itemize}
	
	Die nichttrivialen Elemente können als fundamentale Richtungsmoden interpretiert werden.
	
	\section{Tensorielle Erweiterung}
	
	\subsection{Definition des EABC-Tensors}
	
	Wir definieren den allgemeinen EABC-Tensor durch
	
	\[
	T(a,b,c)=aA+bB+cC.
	\]
	
	In Matrixdarstellung erhält man:
	
	\[
	T=
	\frac1{\sqrt2}
	\begin{pmatrix}
		0&a&b\\
		a&0&c\\
		b&c&0
	\end{pmatrix}.
	\]
	
	\begin{definition}
		Der Raum aller Tensoren der Form
		
		\[
		T(a,b,c)
		\]
		
		heißt \emph{EABC-Tensorraum}.
	\end{definition}
	
	\section{Grundlegende Eigenschaften}
	
	\subsection{Symmetrie}
	
	Der Tensor ist symmetrisch:
	
	\[
	T^T=T.
	\]
	
	Daher besitzt er ausschließlich reelle Eigenwerte.
	
	\subsection{Spurfreiheit}
	
	Es gilt:
	
	\[
	\operatorname{tr}(T)=0.
	\]
	
	Somit beschreibt der Tensor reine Kopplungsstrukturen ohne isotropen Anteil.
	
	\subsection{Kopplungscharakter}
	
	Alle Informationen befinden sich ausschließlich in den Offdiagonalelementen:
	
	\[
	a,b,c.
	\]
	
	Dies bedeutet geometrisch:
	
	\begin{itemize}
		\item keine Selbstkopplung,
		\item ausschließlich Wechselwirkungen,
		\item reine Übergangsstruktur.
	\end{itemize}
	
	\section{Die charakteristische Gleichung}
	
	Die charakteristische Gleichung lautet:
	
	\[
	\det(T-\lambda I)=0.
	\]
	
	Direkte Berechnung ergibt:
	
	\[
	\lambda^3
	-\frac12(a^2+b^2+c^2)\lambda
	-\frac{abc}{\sqrt2}=0.
	\]
	
	Dies ist die fundamentale Spektralgleichung des EABC-Tensorraums.
	
	\section{Die fundamentalen Invarianten}
	
	\subsection{Quadratische Invariante}
	
	Die erste fundamentale Invariante lautet:
	
	\[
	I_2=a^2+b^2+c^2.
	\]
	
	Sie misst die Gesamtstärke der Kopplung.
	
	\subsection{Kubische Invariante}
	
	Die zweite fundamentale Invariante lautet:
	
	\[
	I_3=abc.
	\]
	
	Diese Größe verschwindet sofort, sobald eine Richtung entfällt.
	
	\begin{bemerkung}
		Die kubische Invariante misst somit eine genuine dreidimensionale Kohärenz.
	\end{bemerkung}
	
	\section{Spektraltheorie}
	
	Die charakteristische Gleichung kann geschrieben werden als:
	
	\[
	\lambda^3-p\lambda-q=0
	\]
	
	mit
	
	\[
	p=\frac12I_2,
	\qquad
	q=\frac{I_3}{\sqrt2}.
	\]
	
	\subsection{Diskriminante}
	
	Die Diskriminante der kubischen Gleichung lautet:
	
	\[
	\Delta=4p^3-27q^2.
	\]
	
	Einsetzen ergibt:
	
	\[
	\Delta=
	\frac12(a^2+b^2+c^2)^3
	-\frac{27}{2}a^2b^2c^2.
	\]
	
	\begin{satz}
		Für alle reellen Werte \(a,b,c\) gilt:
		
		\[
		\Delta\ge0.
		\]
	\end{satz}
	
	\begin{proof}
		Nach der Ungleichung vom arithmetischen und geometrischen Mittel gilt:
		
		\[
		(a^2+b^2+c^2)^3
		\ge
		27a^2b^2c^2.
		\]
		
		Einsetzen liefert unmittelbar:
		
		\[
		\Delta\ge0.
		\]
	\end{proof}
	
	\begin{korollar}
		Der EABC-Tensor besitzt stets ausschließlich reelle Eigenwerte.
	\end{korollar}
	
	\section{Quaternionische Interpretation}
	
	Die EABC-Struktur besitzt eine natürliche Quaternionenabbildung:
	
	\[
	A\leftrightarrow i,
	\qquad
	B\leftrightarrow j,
	\qquad
	C\leftrightarrow k.
	\]
	
	Der Tensor entspricht daher einem quaternionischen Richtungsvektor:
	
	\[
	q=ai+bj+ck.
	\]
	
	Die quadratische Invariante entspricht exakt der Quaternionennorm:
	
	\[
	|q|^2=a^2+b^2+c^2.
	\]
	
	\section{Geometrische Interpretation}
	
	\subsection{Tensor als Kopplungsraum}
	
	Die Größen
	
	\[
	a,b,c
	\]
	
	sind keine gewöhnlichen Koordinaten, sondern Kopplungsintensitäten zwischen fundamentalen Richtungsmoden.
	
	\subsection{Orientierte Triple-Kohärenz}
	
	Die kubische Invariante
	
	\[
	abc
	\]
	
	wirkt wie ein orientiertes Volumenelement.
	
	Sie beschreibt:
	
	\begin{itemize}
		\item chirale Struktur,
		\item dreifache Resonanz,
		\item orientierte Kohärenz,
		\item geometrische Nichtseparierbarkeit.
	\end{itemize}
	
	\section{Spezialfall maximaler Symmetrie}
	
	Für
	
	\[
	a=b=c=\alpha
	\]
	
	erhält man:
	
	\[
	T=
	\frac{\alpha}{\sqrt2}
	\begin{pmatrix}
		0&1&1\\
		1&0&1\\
		1&1&0
	\end{pmatrix}.
	\]
	
	Die Eigenwerte lauten:
	
	\[
	\lambda_1=\sqrt2\alpha
	\]
	
	und
	
	\[
	\lambda_2=\lambda_3=-\frac{\alpha}{\sqrt2}.
	\]
	
	\begin{bemerkung}
		Dies entspricht einem dominanten kollektiven Modus mit zweifach entarteter transversaler Struktur.
	\end{bemerkung}
	
	\section{Physikalische Interpretation}
	
	\subsection{Kopplungstensor}
	
	Der Tensor kann als allgemeiner Wechselwirkungsoperator interpretiert werden.
	
	\subsection{Anisotroper Spannungstensor}
	
	Da die Diagonale verschwindet, beschreibt der Tensor reine Scherungs- und Kopplungsanteile.
	
	\subsection{Spinartige Struktur}
	
	Die spurfrei-symmetrische Form erinnert an Spin-2-Felder und Gravitationstensoren.
	
	\section{Verbindung zur Random-Matrix-Theorie}
	
	Der Tensor besitzt strukturelle Ähnlichkeiten zu Wigner-Matrizen:
	
	\begin{itemize}
		\item Symmetrie,
		\item Offdiagonalkopplung,
		\item reelles Spektrum,
		\item Resonanzstruktur.
	\end{itemize}
	
	Im Gegensatz zur Random-Matrix-Theorie handelt es sich hier jedoch um ein vollständig deterministisches System.
	
	\section{EABC-Spektralgeometrie}
	
	Die vollständige Spektralstruktur wird allein durch zwei Größen bestimmt:
	
	\[
	I_2
	\quad\text{und}\quad
	I_3.
	\]
	
	Dies macht den EABC-Tensorraum zu einem minimalen spektralen Modell geometrischer Kopplung.
	
	\section{Ausblick}
	
	Die vorgestellte Tensorstruktur eröffnet zahlreiche mögliche Erweiterungen:
	
	\begin{itemize}
		\item höherdimensionale EABC-Tensoren,
		\item quaternionische Holonomien,
		\item spektrale Netzwerke,
		\item topologische Phasenräume,
		\item gravitationsartige Tensorfelder,
		\item Verbindung zu SU(2)- und SO(3)-Darstellungen,
		\item Anwendungen in Zahlentheorie und Primzahlspektren.
	\end{itemize}
	
	Insbesondere könnte die kubische Invariante eine fundamentale Rolle in orientierten diskreten Geometrien spielen.
	
	\section{Schlussbemerkung}
	
	Der EABC-Tensorraum stellt ein bemerkenswert einfaches, aber strukturell tiefes mathematisches Objekt dar.
	
	Trotz minimaler algebraischer Voraussetzungen entstehen:
	
	\begin{itemize}
		\item nichttriviale Spektralstrukturen,
		\item stabile reelle Eigenwertsysteme,
		\item quaternionische Geometrien,
		\item orientierte Kohärenzinvarianten,
		\item tensorielle Resonanzräume.
	\end{itemize}
	
	Dies deutet darauf hin, dass die EABC-Struktur als fundamentale geometrische Kopplungsalgebra interpretiert werden kann.
	
	\vspace{1cm}
	
	\begin{center}
		\textit{„Nicht die Punkte tragen die Struktur,\\
			sondern ihre Wechselwirkungen.“}
	\end{center}
	
\end{document}